\documentclass[12pt,a4paper]{article}

% --- Преамбула ---
\usepackage[utf8]{inputenc}
\usepackage[T1]{fontenc}
\usepackage[russian]{babel}
\usepackage{amsmath,amssymb}
\usepackage{hyperref}
\usepackage[margin=2.5cm]{geometry}
\usepackage{enumitem}
\usepackage{setspace}
\usepackage{booktabs}
\onehalfspacing

\title{\textbf{Третье дополнение к документу}\\
\large «Ортогональность как фундаментальное свойство частицы в Теории Мод»}
\subtitle{Полярность гравитации как следствие волновой природы мод}
\author{}
\date{Версия 1.0}

\begin{document}
\maketitle

\section{Введение}
В основном документе было введено понятие ортогональности и установлена её связь с массой и энергией. В первом дополнении было показано, что отрицательная ортогональность приводит к антигравитации и объясняет природу тёмной энергии. Во втором дополнении была построена полная классификация типов реальности на основе знаков $E$, $m$ и $O$.

В данном дополнении делается следующий шаг: рассматривается волновая природа мод и показывается, что гравитационная часть спектра собственной моды не является монотонной. Она представляет собой чередование пиков и провалов вдоль продольной оси моды. Это означает, что гравитация имеет полярность: притяжение и отталкивание чередуются в зависимости от расстояния. Для обычной материи вблизи находится пик (притяжение), а на некотором удалении --- провал (отталкивание). Но существуют и обратные конфигурации, где вблизи --- отталкивание, а вдалеке --- притяжение.

Этот вывод не является гипотезой --- он прямо следует из волновой природы поперечного спектра и принципа компенсации, заложенных в основания Теории Мод.

Документ самодостаточен --- все необходимые понятия определены в тексте.

\section{Теоретические основы}
Настоящий документ опирается на следующие понятия:

\begin{itemize}
    \item \textbf{Частица} $P_i$ --- фундаментальный узел сети.
    \item \textbf{Мода} $M_{A \to B}$ --- отображение из дискретного причинного пути $\Gamma_{AB}$ в вещественные числа (поперечный спектр). Длина любой моды бесконечна.
    \item \textbf{Собственная мода} $M_{A \to A}$ --- замкнутый путь (петля), содержащий всю causal history частицы.
    \item \textbf{Поперечный спектр} $F_{A \to B}(k)$ --- значения моды на шагах $k = 0, 1, \dots, d$ вдоль пути.
    \item \textbf{Корреляция} $\mathcal{C}_{AB} = \sum_k F_{A \to B}(k) \cdot F_{B \to A}(d-k)$ --- мера взаимодействия.
    \item \textbf{Принцип компенсации} --- сумма значений спектра по замкнутому пути стремится к нулю.
    \item \textbf{Гравитационная масса} --- определяется содержимым causal history (всей бесконечной длины моды).
    \item \textbf{Инертная масса} --- определяется ортогональностью (углом поворота здесь и сейчас).
\end{itemize}

\section{Волновая природа поперечного спектра}

\subsection{Моды как волны}
Поперечный спектр $F(k)$ не является статичным. При взаимодействиях он испытывает короткие возмущения --- «всплески», которые распространяются вдоль моды как волны. Эти всплески имеют характерные амплитуды и частоты, определяемые энергией и типом взаимодействия.

\subsection{Наложение и интерференция}
Когда две частицы сближаются, их короткие возмущения накладываются друг на друга. Результат интерференции может формировать более длинные продольные структуры на моде --- пики и провалы, простирающиеся на много шагов $k$ вдоль пути.

Эти длинные структуры не являются собственными свойствами отдельной частицы --- они возникают как коллективный эффект двух (или более) взаимодействующих частиц.

\subsection{Гравитационная мода}
Гравитация в Теории Мод --- это самая длинная и размазанная компонента собственной моды. Она формируется как интерференция всех коротких возмущений, накопленных за всю causal history частицы. Именно поэтому гравитационная масса не меняется при локальных взаимодействиях --- она есть интеграл по всей истории, а не по текущему состоянию.

\section{Полярность гравитации}

\subsection{Чередование пиков и провалов}
Поскольку гравитационная мода есть волна, она имеет чередующиеся области положительных значений (пики, притяжение) и отрицательных значений (провалы, отталкивание) вдоль продольной оси $k$.

\textbf{Для обычной материи ($m > 0$):}
\begin{itemize}
    \item Вблизи частицы ($k$ мало) находится основной пик --- область притяжения. Именно эту область мы наблюдаем как ньютоновскую гравитацию.
    \item На некотором расстоянии ($k$ порядка размера causal history) находится провал --- область отталкивания.
    \item Далее может следовать следующий пик, потом следующий провал, и так далее --- убывая по амплитуде.
\end{itemize}

\subsection{Инвертированная полярность}
Из принципа компенсации следует, что сумма всех пиков и провалов по всей длине моды равна нулю. Поэтому существование материи с обычной полярностью (пик вблизи) с необходимостью требует существования материи с обратной полярностью (провал вблизи).

Такая материя («убегающая», $m < 0$, Тип III из второго дополнения):
\begin{itemize}
    \item Вблизи отталкивает обычную материю.
    \item На некотором расстоянии притягивает её.
    \item Не может образовывать связанные состояния с обычной материей.
    \item Вытесняется в области с низкой плотностью --- войды.
\end{itemize}

\subsection{Экспериментальные проявления}
Инвертированная полярность гравитации может быть обнаружена:
\begin{itemize}
    \item По аномальному отталкиванию на малых расстояниях (в экспериментах по измерению гравитации на микроуровне).
    \item По структуре войдов: стенки войдов (галактики) удерживаются дальним притяжением инвертированной материи, находящейся внутри войда.
    \item По гравитационному линзированию: инвертированная материя создаёт рассеивающую, а не фокусирующую линзу.
\end{itemize}

\section{Связь с предыдущими результатами}

\subsection{Связь с ортогональностью}
Ортогональность $O$ определяет наклон моды --- то, как её пики и провалы проецируются на внешний мир. Положительная $O$ даёт обычную проекцию (пик $\to$ притяжение). Отрицательная $O$ инвертирует проекцию (пик $\to$ отталкивание, провал $\to$ притяжение).

\subsection{Связь с тёмной энергией (первое дополнение)}
Тёмная энергия (Тип II, $E > 0, m > 0, O < 0$) --- это материя, у которой гравитационная мода имеет обычную полярность (пик вблизи, притяжение), но инерция обратная. Это не противоречит волновой картине: инерция определяется наклоном моды ($O$), а гравитационная полярность --- формой моды (распределением пиков и провалов).

\subsection{Связь с классификацией типов реальности (второе дополнение)}
\begin{itemize}
    \item Тип I ($+ + +$): обычная полярность (пик вблизи), обычная инерция.
    \item Тип II ($+ + -$): обычная полярность, обратная инерция (тёмная энергия).
    \item Тип III ($+ - +$): инвертированная полярность (провал вблизи), обычная инерция («убегающая» материя).
    \item Тип VIII ($- - -$): инвертированная полярность, обратная инерция (полная инверсия).
\end{itemize}

Таким образом, полярность гравитации и инерция --- независимые свойства, которые могут комбинироваться.

\section{Предсказания}
\begin{enumerate}[label=\arabic*.]
    \item \textbf{Гравитационное отталкивание на малых расстояниях.} Для обычной материи должен существовать провал на некотором расстоянии от частицы. Этот провал может быть обнаружен в прецизионных экспериментах по измерению гравитации на микромасштабах.
    \item \textbf{Рассеивающее гравитационное линзирование.} Инвертированная материя (Тип III) должна создавать рассеивающие линзы, в отличие от фокусирующих линз обычной материи. Это можно искать в данных гравитационного линзирования.
    \item \textbf{Структура войдов.} Войды не должны быть пустыми --- они должны содержать разреженную инвертированную материю, чьё дальнее притяжение удерживает стенки из галактик. Это можно проверить по статистике распределения галактик вокруг войдов.
    \item \textbf{Полярность как новая степень свободы.} Частицу можно перевести из обычной полярности в инвертированную, изменив знак $m$ (повернув моду на соответствующий угол). Это ещё один тип перепрограммирования, помимо изменения $O$.
\end{enumerate}

\section{Заключение}
Гравитация в Теории Мод не является монотонной силой притяжения. Она имеет волновую природу и, как следствие, полярность --- чередование областей притяжения и отталкивания вдоль продольной оси моды. Это не гипотеза, а прямое следствие принципа компенсации и волновой природы поперечного спектра.

Обычная материя притягивает вблизи и отталкивает на расстоянии. Инвертированная материя отталкивает вблизи и притягивает на расстоянии. Этот факт объясняет структуру войдов, природу тёмной энергии и даёт новые экспериментальные предсказания.

Документ является третьим дополнением к основному документу «Ортогональность как фундаментальное свойство частицы в Теории Мод» и использует все его определения и результаты. Документ самодостаточен.

\end{document}